\section{Introduction}
\label{sec:introduction}

\subsection{The Local-Optimum Problem in Single-Chain MCMC}

The B-series redistricting platform constructs congressional maps via a
bisection tree: the state graph is recursively split into $k$ population-balanced
districts by minimising the weighted edge cut at each node~\citep{b8paper}.
The Search layer determines how many MCMC steps are taken and which resulting
plan is returned.
For most states with small to medium district counts ($k \le 20$), a single
\fr{} chain run for $T=1{,}000$ steps adequately explores the feasible plan
space: the chain mixes freely within its balance-tolerance constraint and
the minimum edge-cut plan is selected by percentile.

For large-$k$ states, this single-chain strategy fails in a characteristic way.
Texas ($k=38$), California ($k=52$), and Florida ($k=28$) have state graphs
with tens of thousands of census tracts.
The feasible plan space at strict balance tolerance ($\varepsilon = 0.5\%$) has
a rugged energy landscape: many local-optimum plans separated by high edge-cut
``barriers.''
The single \fr{} chain gets trapped in one basin of attraction; running for
more steps explores the basin thoroughly but cannot escape to qualitatively
different plan regions without passing through infeasible or high-EC intermediate
states.

\paragraph{Track 1: Geographic coarsening.}
\ms{} (G.11)~\citep{dellaLibera2026multiscale} addresses this by operating at
a coarser spatial resolution: groups of tracts are merged into super-tracts,
and the chain proposes swaps at the super-tract level.
This enlarges the proposal neighbourhood and allows larger structural changes
per step.
The cost is that coarsening introduces approximation; the final plan is
projected back to tract resolution.

\paragraph{Track 2: Temperature ladders.}
This paper introduces a second track: \textbf{Parallel Tempering} (PT), which
runs $N_{\mathrm{replicas}}$ chains simultaneously, each at a different balance
tolerance.
Hot chains (high $\varepsilon$) mix freely because they accept almost any
balanced plan; cold chains (low $\varepsilon$) produce distribution-correct
samples but mix slowly.
After every $\mathtt{swap\_interval}$ steps, adjacent chains propose to exchange
their current plans via a Metropolis criterion.
Accepted swaps pass global diversity downward from the hot chains to the cold
chain.

\subsection{Parallel Tempering as a Redistricting Search Strategy}

In statistical mechanics, parallel tempering (PT) is a method for sampling from
complex energy landscapes~\citep{swendsen1986, geyer1991}.
The ``temperature'' in statistical mechanics controls how strongly a system
is constrained to low-energy states; high temperature permits wide exploration
of state space.
In redistricting, the analogous parameter is the population balance tolerance
$\varepsilon$: a high tolerance accepts more plans as feasible, giving the chain
a much larger proposal neighbourhood and faster mixing.

The key insight is that $\varepsilon$ need not be fixed for all chains.
In \pt{} redistricting:
\begin{itemize}
\item The \emph{cold chain} ($\varepsilon_0 = 0.5\%$) targets the statutory
  distribution.
  Only its samples are used for plan selection.
\item The \emph{hot chain} ($\varepsilon_{N-1} = 5\%$) mixes rapidly because
  it can move between any roughly population-balanced plans.
  It serves only as a diversity source.
\item Intermediate chains form a geometric ladder between cold and hot.
\end{itemize}
When a hot chain reaches a low-EC plan that is far from the cold chain's current
location in plan space, a swap can transfer this plan to the cold chain, which
then explores the new basin under tight balance constraints.
This is the mechanism by which PT escapes local optima without modifying the
distributional target of the cold chain.

\subsection{Main Results}

\paragraph{Algorithm.}
We define \pt{} redistricting (Section~\ref{sec:algorithm}) as a Search-layer
strategy that runs $N_{\mathrm{replicas}}$ \fr{} chains on a geometric tolerance
ladder with swap proposals at every $\mathtt{swap\_interval}$ steps.
The cold chain records all visited plans; the plan at the $p$-th percentile of
the cold chain's edge-cut distribution is returned.
The algorithm has four parameters:
\texttt{n\_replicas} (default 4),
\texttt{swap\_interval} (default 10),
\texttt{cold\_tol} (default 0.005),
\texttt{hot\_tol} (default 0.05).

\paragraph{Empirical.}
On TX ($k=38$), \pt{} with $N=4$ achieves lower cold-chain minimum edge cut
than single-chain \fr{}, demonstrating that replica exchange escapes the local
optima where the single chain stalls.
On NC ($k=14$) and WI ($k=8$), the improvement is modest---the single chain
already mixes adequately for these states.
Observed swap acceptance rates of $18$--$24\%$ match the empirical target of
${\approx}23\%$~\citep{kone2005}, providing evidence that the geometric ladder
is well-calibrated for redistricting chains.

\paragraph{Decision rule.}
\pt{} is the recommended search strategy when $k \ge 30$ and local-optimum
trapping is suspected.
For medium $k$ ($8 \le k \le 20$), single-chain \fr{} is preferred for its
lower cost and equivalent quality.
For large $k$ with strong geographic structure, \ms{}~\citep{dellaLibera2026multiscale}
may provide better efficiency per replica.

\subsection{Paper Structure}

Section~\ref{sec:background} reviews parallel tempering theory and the
redistricting context.
Section~\ref{sec:algorithm} gives the formal algorithm, pseudocode, seeding
specification, and audit chain.
Section~\ref{sec:empirical} presents comparisons on NC, WI, and TX.
Section~\ref{sec:when-to-use} provides a structured decision guide.
Section~\ref{sec:conclusion} situates \pt{} in the broader B/G-series algorithm
landscape.

\subsection{Relation to Prior Work}

\pt{} is a Search-layer algorithm, orthogonal to the Structure-layer algorithms
(\geosec{}, \areasec{}, \sa{}) and the bisection-tree SeedCompositor strategies
(\cs{}, \be{}).
\citet{swendsen1986} introduced replica exchange Monte Carlo for spin glasses;
\citet{geyer1991} established the theoretical basis for using the cold chain
output as a valid sample from the target distribution.
\citet{kone2005} derived the ${\approx}23\%$ optimal swap acceptance rate for
Gaussian continuous systems; we adopt this as an empirical calibration target
for discrete redistricting chains, not as a derived guarantee.
\ms{}~\citep{dellaLibera2026multiscale} is the closest B/G-series alternative,
targeting the same large-$k$ mixing problem via geographic coarsening
rather than temperature ladders.
\pt{} differs in mechanism: it requires no coarsening and preserves the exact
\fr{} proposal distribution at every tolerance level.
